\section{Airborne vs.\ Mechanical Coupling}\label{sec:coupling}

This section presents the central novel result of this work: a quantitative
comparison of airborne acoustic and mechanical (whole-body vibration) coupling
to the abdominal flexural modes derived in \S\ref{sec:formulation}.

\subsection{Airborne Acoustic Pathway}

Table~\ref{tab:airborne} summarises the predicted airborne displacement at
resonance for the $n=2$ flexural mode across a range of sound pressure levels.

\begin{table}[htbp]
  \centering
  \caption{Airborne displacement at $n=2$ resonance ($f_2 = \SI{4.4}{\hertz}$,
  $Q = 3.3$, $E = \SI{0.1}{\mega\pascal}$).}
  \label{tab:airborne}
  \begin{tabular}{rcccc}
    \hline
    SPL (dB) & $p_\mathrm{inc}$ (Pa) & $p_\mathrm{eff}$ (Pa) &
    $\xi$ ($\mu$m) & PIEZO? \\
    \hline
    100 & 2.0 & $3.2 \times 10^{-4}$ & 0.014 & No \\
    110 & 6.3 & $1.0 \times 10^{-3}$ & 0.045 & No \\
    120 & 20.0 & $3.2 \times 10^{-3}$ & 0.14  & No \\
    130 & 63.2 & $1.0 \times 10^{-2}$ & 0.45  & No \\
    140 & 200.0 & $3.2 \times 10^{-2}$ & 1.44  & Marginal \\
    \hline
  \end{tabular}
\end{table}

At \SI{120}{\decibel} SPL (the threshold of pain), the predicted wall
displacement is \SI{0.14}{\micro\metre}---well below the $0.5$--$\SI{2.0}
{\micro\metre}$ activation range of mechanosensitive PIEZO channels.  Reaching
the \SI{1}{\micro\metre} threshold requires approximately \SI{137}{\decibel},
a level associated with close proximity to jet engines and acoustic weapons.

\subsection{Mechanical (Whole-Body Vibration) Pathway}

In contrast, mechanical vibration transmitted through the skeleton couples
directly to the shell without the $(ka)^n$ penalty.
Table~\ref{tab:mechanical} shows the predicted relative displacement for the
$n=2$ mode under whole-body vibration.

\begin{table}[htbp]
  \centering
  \caption{Mechanical coupling at $n=2$ resonance
  ($f_2 = \SI{4.4}{\hertz}$, $Q = 3.3$).}
  \label{tab:mechanical}
  \begin{tabular}{rccc}
    \hline
    $a_\mathrm{rms}$ (m/s$^2$) & $x_\mathrm{base}$ ($\mu$m) &
    $\xi_\mathrm{rel}$ ($\mu$m) & PIEZO? \\
    \hline
    0.1 & 187 & 623 & Yes \\
    0.5 (EU action) & 935 & 3115 & Yes \\
    1.0 & 1869 & 6230 & Yes \\
    1.15 (EU limit) & 2150 & 7165 & Yes \\
    \hline
  \end{tabular}
\end{table}

At the EU Whole-Body Vibration Directive action value of
\SI{0.5}{\metre\per\second\squared}, the predicted displacement is
\SI{3115}{\micro\metre}---exceeding the PIEZO threshold by more than three
orders of magnitude.

\subsection{Coupling Ratio}

The ratio of mechanical to airborne displacement defines the coupling
disparity:
\begin{equation}
  \mathcal{R} = \frac{\xi_\mathrm{mech}}{\xi_\mathrm{air}}
  \sim 10^3 \text{--} 10^4.
\end{equation}
This ratio depends primarily on $(ka)^{-n}$ and the base-excitation
amplification $H_\mathrm{rel}$, and is robust to uncertainties in material
properties because both pathways excite the same modes.

\subsection{Energy Budget}

To verify self-consistency, we compute the absorption cross-section using
reciprocity:
\begin{equation}
  \sigma_\mathrm{abs} = \frac{(2n+1)\lambda^2}{4\pi}
  \,\frac{\Gamma_\mathrm{rad}}{\Gamma_\mathrm{total}},
\end{equation}
where $\Gamma_\mathrm{rad}$ is the radiation damping width and
$\Gamma_\mathrm{total} = \Gamma_\mathrm{rad} + \Gamma_\mathrm{struct}$.
For the $n=2$ mode in air, $\zeta_\mathrm{rad} \approx 10^{-15}$ whereas
$\zeta_\mathrm{struct} = 0.15$, giving an absorption efficiency of order
$10^{-14}$.  The shell is effectively transparent to airborne infrasound.

\subsection{Comparison with Empirical Evidence}

The model predictions are consistent with two independent bodies of evidence:
\begin{itemize}
  \item \textbf{WBV studies}: ISO~2631 documents gastrointestinal symptoms at
    $4$--$\SI{8}{\hertz}$ under mechanical vibration \cite{iso2631}.
  \item \textbf{Airborne infrasound studies}: No gastrointestinal effects
    have been observed at SPL $\leq \SI{150}{\decibel}$ \cite{Leventhall2007}.
\end{itemize}
The coupling disparity provides a quantitative explanation for this apparent
contradiction: the two bodies of literature address different excitation
mechanisms with fundamentally different coupling efficiencies.
